%==============================================================================
% RESULTS SECTION
%==============================================================================
\section{Results}

\subsection{PDB Structure Validation}

To independently validate the CCD mapper accuracy, we compared GlycoSMILES2BAP output against known correct structures from the Protein Data Bank. Four PDB structures were selected to test critical stereochemistry handling capabilities:

\begin{table}[h]
\centering
\caption{PDB Validation Test Cases}
\begin{tabular}{lllp{5cm}}
\toprule
ID & PDB & Structure & Key Validation Point \\
\midrule
V001 & 5NSC & Fucosylated & FUC = alpha-L-fucose (L-configuration) \\
V002 & 1L6R & Lactose & GAL C4 axial hydroxyl (epimer distinction) \\
V003 & 2VXR & Sialyllactose & SIA anomeric = C2 (ketose handling) \\
V004 & 5K65 & M3 N-glycan & Branch topology preservation \\
\bottomrule
\end{tabular}
\label{tab:pdb_cases}
\end{table}

\subsubsection{Validation Results}

GlycoSMILES2BAP achieved 100\% accuracy on all four test cases:

\begin{table}[h]
\centering
\caption{PDB Validation Results}
\begin{tabular}{llccc}
\toprule
Test Case & Expected CCD & Actual CCD & Anomeric & Status \\
\midrule
Fucosylated (5NSC) & FUC & FUC & C1 & PASS \\
Lactose (1L6R) & GAL & GAL & C1 & PASS \\
Sialyllactose (2VXR) & SIA & SIA & C2 & PASS \\
M3 N-glycan (5K65) & MAN/BMA/NAG & MAN/BMA/NAG & C1 & PASS \\
\bottomrule
\end{tabular}
\label{tab:pdb_results}
\end{table}

\textbf{Critical Finding: Sialic Acid C2 Anomeric Position}

The V003 test case (PDB:2VXR, sialyllactose) represents the most critical validation. Sialic acids (Neu5Ac) are ketoses with the anomeric carbon at C2, not C1. The CCD mapper correctly:
\begin{enumerate}
\item Mapped Neu5Ac to SIA (alpha-D configuration)
\item Assigned anomeric position to C2 (correct for ketose)
\item Assigned ring oxygen to O6 (6-membered ring)
\end{enumerate}

This directly addresses the stereochemistry error identified by Huang et al.\ (2025).

\subsection{Failure Case Analysis}

On the GlyTouCan database processing (n=100), 6 structures failed (94\% success rate). Failure analysis:

\begin{table}[h]
\centering
\caption{Failure Case Analysis}
\begin{tabular}{lll}
\toprule
Failure Type & Count & Description \\
\midrule
Unsupported CCD & 3 & Rare monosaccharides (KDN, unusual modifications) \\
Unusual linkage & 2 & Non-standard positions (1-2, 1-6 branch) \\
Input error & 1 & Malformed IUPAC notation \\
\bottomrule
\end{tabular}
\label{tab:failures}
\end{table}